% The paper's own proof of Theorem 16, in Lean: a readable account.
% Build: lualatex proof-report.tex (twice)
% The Lean statements in this file are copied from the sources of the package.
\documentclass[11pt,a4paper]{article}
\usepackage{fontspec}
\setmonofont{DejaVu Sans Mono}[Scale=0.82]
\usepackage{amsmath,amssymb,amsthm}
\usepackage[margin=2.4cm]{geometry}
\usepackage{booktabs,longtable,array}
\usepackage{fancyvrb}
\usepackage{xcolor}
\usepackage[colorlinks=true,linkcolor=blue!50!black,urlcolor=blue!50!black,citecolor=blue!50!black]{hyperref}
\usepackage{enumitem}
\setlist{itemsep=2pt,topsep=4pt}

\newcommand{\leanus}{\textunderscore\allowbreak}
\newcommand{\lean}[1]{{\ttfamily\let\_\leanus #1}}
\emergencystretch=3em
\newcommand{\GT}{G_{\mathrm{T}}}
\newcommand{\GF}{G_{\mathrm{F}}}
\newcommand{\GB}{G_{\mathrm{B}}}
\newcommand{\T}{\mathrm{T}}
\newcommand{\F}{\mathrm{F}}
\newcommand{\B}{\mathrm{B}}
\DefineVerbatimEnvironment{code}{Verbatim}{fontsize=\footnotesize,frame=leftline,framerule=0.6pt,rulecolor=\color{gray!50},xleftmargin=1em}

\title{The paper's own proof of Theorem~16, in Lean\\[2pt]
\large The statements, the lemma chain, and what is trusted}
\author{Companion to\\ \emph{The Trident: A Minimum Symmetric Clause Gadget for 3SAT to 3-Coloring}\\
M.~Cho, H.~Kim, and I.~Shin}
\date{November 2026}

\begin{document}
\maketitle

\begin{abstract}
\noindent
This note accompanies the Lean~4 package
\lean{own-proof-trident\_minimum\_unique}. The package proves Theorem~16 of
the paper for the paper's Definition~1: every valid $S_3$-symmetric clause
gadget has at least six auxiliary vertices, and one with six is the trident.
The proof follows Section~IV of the paper. Each numbered statement of
Sections~IV-A to IV-C is stated and proved in Lean, and the proof of
Theorem~16 uses these statements at the places where the paper does. There is no search and no
compiled code: the two final theorems depend on Lean's three standard axioms
only, and the package builds in less than a minute once Mathlib is in place.
The note gives the Lean statements next to the statements of the paper,
describes how each proof goes, records what the formalization showed about
the argument, and lists the trusted base and what is \emph{not} formalized.
The formal verification was written with Claude (Anthropic) and is under
evaluation by the authors; it is not part of the article.
\end{abstract}

\tableofcontents

% ---------------------------------------------------------------------------
\section{The claim}\label{sec:claim}

The paper's Definition~1 calls a graph $G$ an \emph{$S_3$-symmetric clause
gadget} for the clause $C=(x_0\lor x_1\lor x_2)$ if it consists of three
pairwise non-adjacent literal vertices $x_0,x_1,x_2$, a palette triangle
$\GT,\GF,\GB$ with $\GB$ adjacent to every literal, and auxiliary vertices,
together with a group action $\rho:S_3\to\operatorname{Aut}(G)$ with
$\rho(\sigma)(x_i)=x_{\sigma(i)}$ that fixes each anchor. The gadget is
\emph{valid} if, with the anchors colored $\T,\F,\B$, a proper 3-coloring
extending a literal coloring in $\{\T,\F\}^3$ exists exactly for the
assignments that satisfy $C$. So every gadget contains six fixed vertices,
the literals and the anchors, and Theorem~16 counts the auxiliary vertices
added to them.

\paragraph{Theorem~16 (paper).} Every valid $S_3$-symmetric clause gadget has
at least six auxiliary vertices, and with six it is the trident, up to a graph
isomorphism that fixes each anchor and each literal vertex.

\medskip
The package proves this statement for gadgets on a finite vertex set
(\lean{Def1.minimum\_unique}), and it proves that the trident is such a
gadget (\lean{Def1.trident\_is\_valid\_gadget}).

\paragraph{Two proofs of one statement.} The package
\lean{compact-proof-trident\_minimum\_unique} proves the same two theorems.
The file \lean{Definition1.lean} is the same in both packages, byte for byte.
The proofs differ. The compact package turns a gadget into a candidate of a
search engine and uses an exhaustive search, seven runs of which are
compiled code. This package uses the lemma chain of the paper and no search.
The two packages share three files: the statement
(\lean{Definition1.lean}), the orbit types of Lemma~6
(\lean{OrbitStructure.lean}) and the validity of the trident
(\lean{TridentGadget.lean}).

% ---------------------------------------------------------------------------
\section{The statement in Lean}\label{sec:statement}

\lean{proof/Definition1.lean} (97 lines) states Definition~1 and the trident
with Mathlib's simple graphs and permutations. The two theorems at the end of
\lean{proof/Theorem16.lean} use these definitions only, so what the package
proves can be checked from this section. All names are in the namespace
\lean{Def1}.

\subsection{Definition 1}
\begin{code}
structure Gadget (V : Type) where
  G : SimpleGraph V
  lit : Fin 3 → V
  gT : V
  gF : V
  gB : V
  lit_inj : Function.Injective lit
  lit_ne_anchor : ∀ i, lit i ≠ gT ∧ lit i ≠ gF ∧ lit i ≠ gB
  lit_nonadj : ∀ i j, ¬ G.Adj (lit i) (lit j)
  palette : G.Adj gT gF ∧ G.Adj gF gB ∧ G.Adj gB gT
  fan : ∀ i, G.Adj gB (lit i)
  ρ : Equiv.Perm (Fin 3) →* (G ≃g G)
  ρ_lit : ∀ σ i, ρ σ (lit i) = lit (σ i)
  ρ_anchor : ∀ σ, ρ σ gT = gT ∧ ρ σ gF = gF ∧ ρ σ gB = gB
\end{code}
\begin{itemize}
\item \lean{G} is a simple graph on the vertex type \lean{V}: undirected and
without loops.
\item \lean{lit 0}, \lean{lit 1}, \lean{lit 2} are the literals $x_0,x_1,x_2$,
and \lean{gT}, \lean{gF}, \lean{gB} are the anchors. The six are distinct:
two literals by \lean{lit\_inj}, a literal and an anchor by
\lean{lit\_ne\_anchor}, and two anchors because they are adjacent.
\item \lean{lit\_nonadj}, \lean{palette}, \lean{fan}: the literals are
pairwise non-adjacent, the anchors form a triangle, and $\GB$ is adjacent to
every literal.
\item \lean{ρ} is a group homomorphism from $S_3$, the permutations of
$\{0,1,2\}$, to the automorphisms of the graph (\lean{G ≃g G}, with
composition). It permutes the literals naturally (\lean{ρ\_lit}) and fixes
each anchor (\lean{ρ\_anchor}).
\end{itemize}
Every other vertex is an auxiliary vertex:
\begin{code}
def Gadget.IsAux (g : Gadget V) (v : V) : Prop :=
  (∀ i, v ≠ g.lit i) ∧ v ≠ g.gT ∧ v ≠ g.gF ∧ v ≠ g.gB
\end{code}
Nothing else is assumed. In particular an edge between a literal and $\GT$
or $\GF$ is allowed, as in Definition~1 itself; the paper excludes such edges
afterwards by an argument (Section~\ref{sec:conventions}).

\paragraph{The six fixed vertices.} Every gadget contains the three
literals, which carry the input, and the three anchors, whose colors are
fixed. What varies from gadget to gadget is the set of auxiliary vertices,
the edges other than the palette triangle and the edges from $\GB$ to the
literals, and the action on the auxiliary vertices. The results count the
auxiliary vertices added to the six fixed ones. They say nothing about
arbitrary graphs with at most twelve vertices: a graph without this frame,
or with a symmetry that moves an anchor, is not a gadget of Definition~1.

\begin{code}
def Gadget.Extends (g : Gadget V) (α : Fin 3 → Bool) : Prop :=
  ∃ c : V → Fin 3, (∀ u v, g.G.Adj u v → c u ≠ c v) ∧
    c g.gT = 0 ∧ c g.gF = 1 ∧ c g.gB = 2 ∧
    ∀ i, c (g.lit i) = if α i then 0 else 1

def Gadget.Valid (g : Gadget V) : Prop :=
  ∀ α : Fin 3 → Bool, g.Extends α ↔ ∃ i, α i = true
\end{code}
The colors $0,1,2$ stand for $\T,\F,\B$. \lean{g.Extends α} says that some
proper 3-coloring gives the anchors the colors $\T,\F,\B$ and gives literal
$i$ the color $\T$ if \lean{α i} is true and $\F$ otherwise. \lean{g.Valid}
says that this holds exactly for the inputs \lean{α} with a true literal,
that is, for the inputs that satisfy $C$.

\subsection{The trident}
\begin{code}
inductive TV
  | x (i : Fin 3)
  | p (i : Fin 3)
  | u (i : Fin 3)
  | gT
  | gF
  | gB
  deriving DecidableEq, Fintype

def tridentEdge : TV → TV → Prop
  | .x i, .p j => i = j
  | .p i, .u j => i = j
  | .p _, .gT => True
  | .u i, .u j => i ≠ j
  | .x _, .gB => True
  | .gT, .gF => True
  | .gF, .gB => True
  | .gB, .gT => True
  | _, _ => False

def trident : SimpleGraph TV := SimpleGraph.fromRel tridentEdge
\end{code}
\lean{SimpleGraph.fromRel} makes the relation symmetric and removes loops.
So \lean{trident} has twelve vertices and the 18 edges $x_ip_i$, $p_iu_i$,
$p_i\GT$, $u_iu_j$ ($i\ne j$), $x_i\GB$ and the palette triangle
(Fig.~2(c) of the paper).

\subsection{The two theorems}
\begin{code}
theorem trident_is_valid_gadget :
    (∃ g : Gadget TV, g.G = trident ∧ g.lit = TV.x ∧ g.gT = TV.gT ∧ g.gF = TV.gF ∧
      g.gB = TV.gB ∧ g.Valid ∧ Nat.card {v // g.IsAux v} = 6) ∧ Fintype.card TV = 12

theorem minimum_unique [Fintype V] (g : Gadget V) (hv : g.Valid) :
    6 ≤ Nat.card {v // g.IsAux v} ∧ 12 ≤ Fintype.card V ∧
    (Nat.card {v // g.IsAux v} = 6 → ∃ e : g.G ≃g trident,
      (∀ i, e (g.lit i) = TV.x i) ∧ e g.gT = TV.gT ∧ e g.gF = TV.gF ∧ e g.gB = TV.gB)
\end{code}
The first theorem says that the trident graph, with the literals $x_i$ and
the three anchors, carries an action as in Definition~1 and is valid (the
paper's Theorem~4). It has six auxiliary vertices, the $p_i$ and the $u_i$,
and so twelve vertices in total. The second is Theorem~16. Let \lean{g} be a
valid gadget on a finite vertex type \lean{V}, and let $k$ be the number of
its auxiliary vertices, \lean{Nat.card \{v // g.IsAux v\}}. The theorem says
that $k\ge6$ and that \lean{g} therefore has at least twelve vertices, the
six fixed ones included. If $k=6$, there is a graph isomorphism \lean{e}
from \lean{g.G} onto the trident that sends each literal to the corresponding
literal and each anchor to the corresponding anchor.

\paragraph{What a reader has to accept.} The seven definitions above,
Mathlib's notions of simple graph, permutation, group homomorphism and graph
isomorphism, and the trusted base of Section~\ref{sec:trust}. The
definitions of the following sections do not occur in the two theorems;
they matter only for the lemmas.

% ---------------------------------------------------------------------------
\section{Conventions of the proof}\label{sec:conventions}

\lean{proof/Basics.lean} fixes the vocabulary of the lemmas. In the
statements of this note, \lean{g} is a gadget on a vertex type \lean{V}
unless a statement introduces it, and the names are in the namespace
\lean{Def1.Gadget}. The three anchors form a family indexed by their
colors, and a coloring for an input is a predicate of its own:
\begin{code}
def anchor : Fin 3 → V
  | 0 => g.gT
  | 1 => g.gF
  | 2 => g.gB

def IsColoring (α : Fin 3 → Bool) (c : V → Fin 3) : Prop :=
  (∀ u v, g.G.Adj u v → c u ≠ c v) ∧ c g.gT = 0 ∧ c g.gF = 1 ∧ c g.gB = 2 ∧
    ∀ i, c (g.lit i) = if α i then 0 else 1
\end{code}
By definition \lean{g.Extends α} is \lean{∃ c, g.IsColoring α c}. Three kinds
of inputs occur in the proofs: $(\T,\T,\T)$, $(\F,\F,\F)$, and the input
with exactly one true literal, for example $(\F,\F,\T)$ as
\lean{onlyTrue 2}.
\begin{code}
def allTrue : Fin 3 → Bool := fun _ => true

def allFalse : Fin 3 → Bool := fun _ => false

def onlyTrue (k : Fin 3) : Fin 3 → Bool := fun i => decide (i = k)
\end{code}

\paragraph{The standing assumption.} After Definition~1 the paper shows that
no valid gadget joins a literal to $\GT$ or to $\GF$: by symmetry such an
edge attaches the anchor to all three literals, and then $(\T,\T,\T)$ or
$(\F,\F,\T)$ is rejected. From then on the paper assumes that there is no
such edge.
\begin{code}
def NoLitAnchor : Prop := ∀ i, ¬ g.G.Adj (g.lit i) g.gT ∧ ¬ g.G.Adj (g.lit i) g.gF

theorem Valid.noLitAnchor (hv : g.Valid) : g.NoLitAnchor
\end{code}
A lemma that needs the assumption either has validity among its hypotheses
and derives it, or names it as a hypothesis \lean{h0 : g.NoLitAnchor}.

\paragraph{Counting.} The vertices of a gadget are the three literals, the
three anchors and the auxiliary vertices:
\begin{code}
theorem card_eq [Finite V] : Nat.card V = Nat.card {v // g.IsAux v} + 6
\end{code}

\paragraph{Minimum gadgets.} The paper argues about ``a minimum valid
gadget''. In Lean (\lean{Corollary9.lean}):
\begin{code}
def IsMinimum : Prop :=
  ∀ (V' : Type) [Finite V'] (g' : Gadget V'), g'.Valid →
    Nat.card {v // g.IsAux v} ≤ Nat.card {v // g'.IsAux v}
\end{code}
The vertex types are in \lean{Type}, the lowest universe of types in Lean.
Every finite set has a copy there, so this is no restriction for finite
gadgets. A minimum valid gadget is a valid gadget with this property. It has
at most six auxiliary vertices, since the trident is valid with six
(\lean{IsMinimum.card\_le\_six}).

% ---------------------------------------------------------------------------
\section{Orbit structure (Section~IV-A of the paper)}\label{sec:orbits}

\subsection{Lemma 6 (orbit structure)}\label{sec:lemma6}
\emph{Paper.} Every orbit of auxiliary vertices has size 1, 2, 3 or 6. A
size-3 orbit carries the natural action, and its induced subgraph is a
triangle or three isolated vertices. A size-2 orbit carries the sign action:
its two vertices are fixed by every even permutation and swapped by every
odd one. A size-6 orbit is free. A vertex of a size-1 or size-2 orbit is
adjacent to all three literals or to none, and a size-3 orbit has one of
four literal attachments: empty, matched, crossed, or both.

\smallskip
A size-3 orbit with the natural action is a \emph{layer}:
\begin{code}
structure Layer (g : Gadget V) where
  pt : Fin 3 → V
  map : ∀ (σ : Equiv.Perm (Fin 3)) (i : Fin 3), g.ρ σ (pt i) = pt (σ i)
  inj : Function.Injective pt
  aux : ∀ i, g.IsAux (pt i)
\end{code}
The attachments of a layer \lean{L} are propositions about its vertex of
index~0. The lemmas after them show that they describe all the edges.
\begin{code}
def matched : Prop := g.G.Adj (L.pt 0) (g.lit 0)

def crossed : Prop := g.G.Adj (L.pt 0) (g.lit 1)

def triangle : Prop := g.G.Adj (L.pt 0) (L.pt 1)

theorem lemma6 (g : Gadget V) {v : V} (hv : g.IsAux v) :
    (Nat.card {w // ∃ σ, g.ρ σ v = w} = 1 ∧ (∀ σ, g.ρ σ v = v) ∧
      ((∀ i, g.G.Adj v (g.lit i)) ∨ ∀ i, ¬ g.G.Adj v (g.lit i))) ∨
    (Nat.card {w // ∃ σ, g.ρ σ v = w} = 2 ∧
      (∃ b, b ≠ v ∧ (∀ σ, Equiv.Perm.sign σ = 1 → g.ρ σ v = v ∧ g.ρ σ b = b) ∧
        ∀ σ, Equiv.Perm.sign σ = -1 → g.ρ σ v = b ∧ g.ρ σ b = v) ∧
      ((∀ i, g.G.Adj v (g.lit i)) ∨ ∀ i, ¬ g.G.Adj v (g.lit i))) ∨
    (Nat.card {w // ∃ σ, g.ρ σ v = w} = 3 ∧ ∃ L : g.Layer, (∃ k, L.pt k = v) ∧
      ((∀ i j, i ≠ j → g.G.Adj (L.pt i) (L.pt j)) ∨ ∀ i j, ¬ g.G.Adj (L.pt i) (L.pt j)) ∧
      ∀ i j, g.G.Adj (L.pt i) (g.lit j) ↔ if i = j then L.matched else L.crossed) ∨
    (Nat.card {w // ∃ σ, g.ρ σ v = w} = 6 ∧ ∀ σ, g.ρ σ v = v ↔ σ = 1)
\end{code}
The orbit of \lean{v} is the type \lean{\{w // ∃ σ, g.ρ σ v = w\}}, and the
four alternatives are the four orbit sizes. In the third one, the last line
says that \lean{L.pt i} is adjacent to \lean{lit i} exactly when the
attachment is matched, and to \lean{lit j} with $j\ne i$ exactly when it is
crossed. The two flags give the four patterns of the paper.

\paragraph{Proof.} \lean{proof/OrbitStructure.lean} has the group theory, as
in the paper. By Lagrange's theorem a stabilizer has order 6, 3, 2 or 1. A
subgroup of order 2 is generated by an involution, which fixes an index; a
subgroup of order 3 has index 2 and is the alternating group; so the
stabilizer is $S_3$, $A_3$, the stabilizer of an index, or trivial
(\lean{subgroup\_cases}). The orbit--stabilizer theorem gives the orbit
sizes and the four actions (\lean{orbit\_structure}). The statements about
edges follow from two facts about three indices: $S_3$ moves an ordered pair
of different indices to any other such pair, and the even permutations move
an index to any other index.

\paragraph{The parts.} The later proofs use Lemma~6 through its parts:
\lean{stab\_cases} (the four stabilizers of a vertex),
\lean{Layer.ofNatural} (the layer of a vertex with the natural action),
\lean{sign\_orbit} (the two vertices of a size-2 orbit),
\lean{adj\_lit\_of\_even\_fixed} (all three literals or none),
\lean{Layer.adj\_self} and \lean{Layer.adj\_lit}. Three more statements of
the same kind describe the other edges of a layer. The paper uses them
without a number.
\begin{itemize}
\item \lean{Layer.adj\_anchor}: an anchor is adjacent to all vertices of a
layer or to none (the flag \lean{L.anch t}).
\item \lean{Layer.adj\_fixed}: the same for a vertex that the even
permutations fix.
\item \lean{Layer.adj\_layer}: the wiring between two layers is empty, a
matching, its complement, or both (the flags \lean{L.wm U} and
\lean{L.wc U}). The paper states this in the proof of Lemma~14.
\end{itemize}

\subsection{Definition 2 and Claim 1 (fibers)}
\emph{Paper.} For a free orbit whose vertices each have exactly one literal
neighbor, the fiber of $x_i$ is the set of vertices of the orbit adjacent to
$x_i$. Each fiber contains exactly two vertices.
\begin{code}
def fiber (v : V) (i : Fin 3) : Set V :=
  {w | (∃ σ, g.ρ σ v = w) ∧ g.G.Adj w (g.lit i)}

theorem fiber_card {v : V} (hfree : ∀ σ, g.ρ σ v = v ↔ σ = 1) {k : Fin 3}
    (hone : ∀ j, g.G.Adj v (g.lit j) ↔ j = k) (i : Fin 3) : Nat.card (g.fiber v i) = 2
\end{code}
The Lean proof counts: the vertex \lean{g.ρ σ v} is adjacent to $x_j$
exactly when $j=\sigma(k)$, and two permutations move $k$ to $i$. The paper
argues with the edge orbit instead. Theorem~16 does not use fibers; in
Lemma~15 the paper names the fibers of the two ends of an internal edge, and
the Lean proof uses the literal neighbors of the two ends directly.

\subsection{Lemma 7 (size-2 orbit reduction)}
\emph{Paper.} Suppose a gadget has exactly one size-2 auxiliary orbit
$\{a,b\}$, and all other auxiliary orbits have size 1 or 3. Then $a$ and $b$
have the same neighbors outside $\{a,b\}$, and the graph admits an
$S_3$-action that fixes $a$ and $b$.
\begin{code}
theorem lemma7 (g : Gadget V) {a b : V} (ha : g.IsAux a) (hb : g.IsAux b)
    (heven : ∀ σ, Equiv.Perm.sign σ = 1 → g.ρ σ a = a ∧ g.ρ σ b = b)
    (hodd : ∀ σ, Equiv.Perm.sign σ = -1 → g.ρ σ a = b ∧ g.ρ σ b = a)
    (hother : ∀ w, g.IsAux w → w ≠ a → w ≠ b →
      (∀ σ, g.ρ σ w = w) ∨ ∃ k : Fin 3, ∀ σ, g.ρ σ w = w ↔ σ k = k) :
    (∀ w, w ≠ a → w ≠ b → (g.G.Adj a w ↔ g.G.Adj b w)) ∧
    ∃ g' : Gadget V, (∀ u w, g'.G.Adj u w ↔ g.G.Adj u w) ∧ g'.lit = g.lit ∧
      g'.gT = g.gT ∧ g'.gF = g.gF ∧ g'.gB = g.gB ∧
      (∀ σ, g'.ρ σ a = a ∧ g'.ρ σ b = b) ∧
      ∀ σ w, w ≠ a → w ≠ b → g'.ρ σ w = g.ρ σ w
\end{code}
\lean{heven} and \lean{hodd} say that $\{a,b\}$ is a size-2 orbit, and
\lean{hother} says that every other auxiliary vertex lies in an orbit of
size 1 or 3. The conclusion has two parts. The first is the equality of the
external neighborhoods. The second gives a gadget \lean{g'} on the same
vertices with the same edges, literals and anchors, whose action fixes $a$
and $b$ and agrees with the action of \lean{g} on every other vertex. The
proof has the two steps of the paper: $b$ is fixed by the even permutations,
so it is adjacent to all vertices of a size-3 orbit or to none; an odd
permutation carries a pair $(a,w)$ to $(b,w')$ with $w'$ in the orbit of
$w$. As in the paper, the new action is $\rho'(\sigma)=\rho(\sigma)$ for even $\sigma$
and $\rho'(\sigma)=\tau\rho(\sigma)$ for odd $\sigma$ (\lean{twist}), where
$\tau$ swaps $a$ and $b$ (\lean{swapIso}); it is a homomorphism because the
sign is one and $\tau$ commutes with every $\rho(\sigma)$.

% ---------------------------------------------------------------------------
\section{General elimination principles (Section~IV-B)}\label{sec:principles}

\subsection{Lemma 8 (color-swap symmetry)}
\emph{Paper.} If neither $\GT$ nor $\GF$ is adjacent to a non-anchor vertex,
then $(\T,\T,\T)$ extends to a proper coloring if and only if $(\F,\F,\F)$
does.
\begin{code}
theorem lemma8
    (hno : ∀ v, (∀ t, v ≠ g.anchor t) → ¬ g.G.Adj g.gT v ∧ ¬ g.G.Adj g.gF v) :
    g.Extends allTrue ↔ g.Extends allFalse
\end{code}
The proof swaps the colors $\T$ and $\F$ on the non-anchor vertices
(\lean{swapColors}).

\subsection{Corollary 9 (anchor normalization)}
\emph{Paper.} In a valid gadget, (i) at least one of $\GT$ and $\GF$ is
adjacent to a non-anchor vertex, and (ii) every auxiliary vertex of a
minimum valid gadget is adjacent to at most one palette anchor.
\begin{code}
theorem cor9_i (hv : g.Valid) :
    ∃ v, (∀ t, v ≠ g.anchor t) ∧ (g.G.Adj g.gT v ∨ g.G.Adj g.gF v)

theorem cor9_ii [Finite V] (hv : g.Valid) (hmin : g.IsMinimum) {v : V} (hvaux : g.IsAux v)
    {s t : Fin 3} (hst : s ≠ t) (hs : g.G.Adj v (g.anchor s)) (ht : g.G.Adj v (g.anchor t)) :
    False
\end{code}
Part (ii) follows the paper. A vertex adjacent to the anchors of colors $s$
and $t$ has the third color $r$ in every coloring, and so has every vertex
of its orbit. Hence no edge joins two vertices of the orbit, and none joins
the orbit to the anchor of color $r$. Merging the orbit into this anchor
gives a gadget with fewer auxiliary vertices that accepts the same inputs.

\paragraph{The merge.} \lean{proof/Removal.lean} builds the smaller gadget.
\begin{code}
structure Removal (g : Gadget V) where
  S : V → Prop
  t : Fin 3
  ex : V → Prop
  S_aux : ∀ v, S v → g.IsAux v
  S_inv : ∀ (σ : Equiv.Perm (Fin 3)) (v : V), S v → S (g.ρ σ v)
  ex_inv : ∀ (σ : Equiv.Perm (Fin 3)) (v : V), ex (g.ρ σ v) ↔ ex v
  ex_anchor : ¬ ex (g.anchor t)
\end{code}
\lean{S} is the set of auxiliary vertices to remove, and \lean{ex} marks the
vertices that are joined to the anchor \lean{t} afterwards. Both sets are
invariant under the action. \lean{R.gadget} is a gadget of Definition~1 on
the vertex type \lean{\{v // ¬ R.S v\}}, with the induced edges, the new
anchor edges, and the restricted action. For Corollary~9(ii), \lean{S} is
the orbit and \lean{ex} marks its neighbors. The same construction with
\lean{ex} empty deletes a set of vertices (Section~\ref{sec:triangle}).
\lean{restrict} and \lean{extend} carry colorings between the two gadgets,
and \lean{card\_aux\_lt} says that the new gadget has fewer auxiliary
vertices.

\subsection{Lemma 10 (mimicry) and Claim 2 (palettes)}
\emph{Paper.} Let $g(\alpha)$ be the set of proper colorings of the
auxiliary vertices when the literals are colored by $\alpha$ and the anchor
colors are fixed. If $g(\alpha)\subseteq g(\F,\F,\F)$ for a satisfying
$\alpha$, the gadget is not valid. The palette of an auxiliary vertex is the
set of colors not used by its literal neighbors; if the palette of every
auxiliary vertex under $\alpha$ is contained in its palette under
$(\F,\F,\F)$, then $g(\alpha)\subseteq g(\F,\F,\F)$.
\begin{code}
def AuxColoring (α : Fin 3 → Bool) (a : V → Fin 3) : Prop :=
  ∃ c, g.IsColoring α c ∧ ∀ v, g.IsAux v → c v = a v

theorem lemma10 {α : Fin 3 → Bool} (hα : ∃ i, α i = true)
    (h : ∀ a, g.AuxColoring α a → g.AuxColoring allFalse a) : ¬ g.Valid

def vertexPalette (α : Fin 3 → Bool) (v : V) : Set (Fin 3) :=
  {d | ∀ i, g.G.Adj v (g.lit i) → d ≠ if α i then 0 else 1}

theorem palette_claim (h0 : g.NoLitAnchor) {α : Fin 3 → Bool}
    (hp : ∀ v, g.IsAux v → g.vertexPalette α v ⊆ g.vertexPalette allFalse v) :
    ∀ a, g.AuxColoring α a → g.AuxColoring allFalse a
\end{code}
\lean{g.AuxColoring α a} says that \lean{a} belongs to $g(\alpha)$: it is
the auxiliary part of a proper coloring for \lean{α}. Claim~2 has the
standing assumption as a hypothesis. The proof recolors the literals with
$\F$, and an edge from a literal to $\GT$ or $\GF$ would not survive this.

\subsection{Lemma 11 (the degree-one channel)}
\emph{Paper.} A valid clause gadget contains an auxiliary vertex that is
adjacent to exactly one literal. Otherwise
$g(\F,\F,\T)\subseteq g(\F,\F,\F)$ and the gadget fails by mimicry.
\begin{code}
theorem lemma11 (hv : g.Valid) :
    ∃ v, g.IsAux v ∧ ∃ k : Fin 3, ∀ j, g.G.Adj v (g.lit j) ↔ j = k

theorem lemma11_containment (h0 : g.NoLitAnchor)
    (hno : ¬ ∃ v, g.IsAux v ∧ ∃ k : Fin 3, ∀ j, g.G.Adj v (g.lit j) ↔ j = k) :
    ∀ a, g.AuxColoring (onlyTrue 2) a → g.AuxColoring allFalse a
\end{code}
The second statement is the second sentence of the lemma, and the first
follows from it by Lemma~10. The containment is proved with Claim~2. The
proof does not split by the literal degree: if a vertex sees $x_2$ and $x_2$
is not its only literal neighbor, it sees a false literal under
$(\F,\F,\T)$ as well, so its palette under $(\F,\F,\T)$ avoids $\F$.

\paragraph{The remark on arbitrary clause gadgets.}\label{sec:arbitrary}
The proof in the paper ends with a remark: the argument does not rely on
$S_3$-symmetry and therefore holds for arbitrary clause gadgets.
\lean{proof/ClauseGadget.lean} proves the lemma in that generality. A
clause gadget is the graph of Definition~1 without the action, with the same
notions of auxiliary vertex and validity. In the statements below,
\lean{C} is a clause gadget, and the names are in the namespace
\lean{Def1.ClauseGadget}:
\begin{code}
structure ClauseGadget (V : Type) where
  G : SimpleGraph V
  lit : Fin 3 → V
  gT : V
  gF : V
  gB : V
  lit_inj : Function.Injective lit
  lit_ne_anchor : ∀ i, lit i ≠ gT ∧ lit i ≠ gF ∧ lit i ≠ gB
  lit_nonadj : ∀ i j, ¬ G.Adj (lit i) (lit j)
  palette : G.Adj gT gF ∧ G.Adj gF gB ∧ G.Adj gB gT
  fan : ∀ i, G.Adj gB (lit i)

theorem lemma11 (hv : C.Valid) :
    ∃ v, C.IsAux v ∧ ∃ k : Fin 3, ∀ j, C.G.Adj v (C.lit j) ↔ j = k
\end{code}
The argument recolors the literals with $\F$, so it needs that no edge
joins a literal to $\GF$. \lean{Valid.noLitAnchor}
(Section~\ref{sec:conventions}) excludes the edges between a literal and
$\GT$ or $\GF$ by the symmetry of the gadget. They are excluded without
symmetry as well:
\begin{code}
theorem noLitAnchor (hv : C.Valid) (i : Fin 3) :
    ¬ C.G.Adj (C.lit i) C.gT ∧ ¬ C.G.Adj (C.lit i) C.gF
\end{code}
An edge from the literal $x_i$ to $\GT$ rejects the input in which only
$x_i$ is true, and an edge from $x_i$ to $\GF$ rejects an input in which
$x_i$ is false and another literal is true. A gadget of Definition~1 is a
clause gadget with the same validity and the same auxiliary vertices
(\lean{Gadget.toClauseGadget\_spec}), so this lemma contains Lemma~11.

\subsection{Corollary 12}
\emph{Paper.} If a triangle module is directly connected to the literals,
the gadget admits no proper 3-coloring under $(\T,\T,\T)$.
\begin{code}
theorem cor12 (L : g.Layer) (htri : L.triangle) (hatt : L.matched ∨ L.crossed) :
    ¬ g.Extends allTrue
\end{code}

% ---------------------------------------------------------------------------
\section{Candidate elimination (Section~IV-C)}\label{sec:elimination}

\subsection{Lemma 13 (matched base extensions)}\label{sec:lemma13}
\emph{Paper.} A gadget built on the matched base, an isolated orbit
$\{p_0,p_1,p_2\}$ matched to the literals, and extended by any number of
singleton orbits is invalid, whatever the edges among the singletons and the
invariant anchor attachments are.
\begin{code}
theorem lemma13 (P : g.Layer) (hm : P.matched) (_hnc : ¬ P.crossed) (hiso : ¬ P.triangle)
    (hsing : ∀ w, g.IsAux w → (∀ i, w ≠ P.pt i) → ∀ σ, g.ρ σ w = w) : ¬ g.Valid
\end{code}
\lean{P} with the three hypotheses \lean{hm}, \lean{\_hnc} and \lean{hiso}
is the matched base, and \lean{hsing} says that every auxiliary vertex
outside \lean{P} is fixed by every permutation. Nothing is assumed about the
edges among these vertices or about the anchors. The proof is the recoloring
of the paper: take a coloring for $(\F,\F,\T)$, recolor \lean{P} with $\B$
if the common neighbors of \lean{P} allow it and with $\T$ otherwise, and
recolor the literals with $\F$.

\subsection{Two layers}
Six auxiliary vertices in two size-3 orbits are described by one structure:
\begin{code}
structure TwoLayers (g : Gadget V) where
  L : g.Layer
  U : g.Layer
  disj : ∀ i j, L.pt i ≠ U.pt j
  all : ∀ w, g.IsAux w → (∃ i, w = L.pt i) ∨ ∃ i, w = U.pt i
\end{code}
The twelve vertices are then indexed by the vertices of the trident
(\lean{fr : TV → V}), and by Lemma~6 every adjacency is one of the flags of
the two layers (\lean{rel}, \lean{adj\_fr}).

\subsection{Lemma 14 (two isolated layers)}
\emph{Paper.} No minimum valid gadget has exactly two size-3 auxiliary
orbits that are both isolated layers.
\begin{code}
theorem lemma14 [Finite V] (hv : g.Valid) (hmin : g.IsMinimum) (L U : g.Layer)
    (hdisj : ∀ i j, L.pt i ≠ U.pt j) (hL : ¬ L.triangle) (hU : ¬ U.triangle) : False
\end{code}
\lean{L} and \lean{U} are two disjoint size-3 orbits. In a minimum gadget
they are all the auxiliary vertices, because the budget is six:
\begin{code}
theorem two_layers_exhaust_budget [Finite V] (hmin : g.IsMinimum) (L U : g.Layer)
    (hdisj : ∀ i j, L.pt i ≠ U.pt j) :
    Nat.card {v // g.IsAux v} = 6 ∧
      ∀ w, g.IsAux w → (∃ i, w = L.pt i) ∨ ∃ i, w = U.pt i
\end{code}
So the two orbits are exactly the auxiliary orbits, as in the paper's
statement, and \lean{TwoLayers.not\_both\_isolated} proves the lemma in the
frame of the previous section. The proof has the cases of the paper. Under $(\F,\F,\F)$ the two layers get
the colors $\T$ and $\B$ in one of the two orders, unless both are attached
to $\GT$ or both to $\GB$. Two layers attached to $\GB$ contradict
Corollary~9(i). For two layers attached to $\GT$: without an edge between
them both get $\B$; a layer without literal neighbors gets $\F$ and the
other $\B$; otherwise, under $(\F,\F,\T)$, two vertices of each layer have
color $\B$, and the wiring joins two of them.

\subsection{The triangle analysis}\label{sec:triangle}
\emph{Paper} (the text after Lemma~14). If the second layer $U$ is a
triangle, it is attached neither to the literals nor to an anchor. The
wiring between the base and $U$ is matched ($d=1$): $d=0$, $d=3$ and $d=2$
are excluded. The branch layer is attached to $\GT$.
\begin{code}
theorem triangle_analysis [Finite V] (T : g.TwoLayers) (hv : g.Valid) (hmin : g.IsMinimum)
    (hPm : T.L.matched) (hPc : ¬ T.L.crossed) (hPiso : ¬ T.L.triangle)
    (hUtri : T.U.triangle) :
    (¬ T.U.matched ∧ ¬ T.U.crossed) ∧ (∀ t, ¬ T.U.anch t) ∧
    (T.L.wm T.U ∧ ¬ T.L.wc T.U) ∧ (T.L.anch 0 ∧ ¬ T.L.anch 1 ∧ ¬ T.L.anch 2)
\end{code}
The conclusion has four parts, in the order of the paper: the triangle has
no literal attachment, it has no anchor, the wiring is matched and not
crossed, and the branch layer is attached to $\GT$ and to no other anchor.
The case $d=0$ is a lemma of its own:
\begin{code}
theorem no_disconnected_layer (T : g.TwoLayers) (hv : g.Valid) (hPm : T.L.matched)
    (hPc : ¬ T.L.crossed) (hPiso : ¬ T.L.triangle)
    (hUlit : ¬ T.U.matched ∧ ¬ T.U.crossed) (hUanch : ∀ t, ¬ T.U.anch t)
    (hw : ¬ T.L.wm T.U ∧ ¬ T.L.wc T.U) : False
\end{code}
A second layer without an edge to the rest is deleted
(Section~\ref{sec:principles}), and what is left is the matched base, which
Lemma~13 excludes. This is the paper's sentence that a disconnected second
layer cannot repair the invalid matched base.

\subsection{Lemma 15 (free orbit exclusion)}
\emph{Paper.} A free orbit cannot be the sole auxiliary orbit of a minimum
valid gadget.
\begin{code}
theorem lemma15 [Finite V] (hv : g.Valid) (hmin : g.IsMinimum) {v : V} (hvaux : g.IsAux v)
    (hsole : ∀ w, g.IsAux w → ∃ σ, g.ρ σ v = w) (_hfree : ∀ σ, g.ρ σ v = v ↔ σ = 1) :
    False
\end{code}
\lean{hsole} says that the orbit of \lean{v} contains every auxiliary
vertex, and \lean{\_hfree} that the orbit is free. The proof follows the
paper: one literal neighbor per vertex (Lemma~11), at most one anchor
(Corollary~9(ii)), which is $\GT$ or $\GF$ (Corollary~9(i)); without an
internal edge the orbit gets $\B$ under $(\F,\F,\F)$; with an internal edge,
$(\T,\T,\T)$ for $\GF$ and an input in which the literal neighbors of the
two ends are false for $\GT$ give two adjacent vertices of color $\B$.

% ---------------------------------------------------------------------------
\section{Theorem 16}\label{sec:thm16}

\begin{code}
theorem minimum_is_trident [Finite V] (g : Gadget V) (hv : g.Valid) (hmin : g.IsMinimum) :
    Nat.card {v // g.IsAux v} = 6 ∧ ∃ e : g.G ≃g trident,
      (∀ i, e (g.lit i) = TV.x i) ∧ e g.gT = TV.gT ∧ e g.gF = TV.gF ∧ e g.gB = TV.gB
\end{code}
This is the proof of the paper for a minimum valid gadget $G$. It has at most
six auxiliary vertices, since the trident is valid with six.

\begin{center}
\begin{tabular}{@{}p{0.2\textwidth}p{0.5\textwidth}p{0.22\textwidth}@{}}
\toprule
step of the paper & what the Lean proof does & uses\\
\midrule
Required first layer & A free orbit has six vertices, so it would be the
sole auxiliary orbit. A vertex $v$ with exactly one literal neighbor exists.
Its orbit has size 3, and its attachment is matched and not crossed. This
orbit is not a triangle. & Lemma~15, Lemma~11, Lemma~6, Corollary~12\\
\addlinespace
Remaining budget & If the other auxiliary vertices do not form a second
size-3 orbit, every other orbit has size 1 or 2. A second size-2 orbit would
need seven auxiliary vertices. One size-2 orbit becomes two singletons. &
Lemma~6, Lemma~7, Lemma~13\\
\addlinespace
Equality case & A second size-3 orbit uses all six vertices. It is not an
isolated layer, so it is a triangle, and the flags of the two layers are
those of the trident. & Lemma~14, the triangle analysis\\
\bottomrule
\end{tabular}
\end{center}

\paragraph{The isomorphism.} \lean{TwoLayers.iso\_trident} turns the flags
into the isomorphism. The indexing \lean{fr : TV → V} is a bijection, and
with the flags left by the triangle analysis the relation \lean{rel} is the
adjacency of the trident; this is checked for the 36 pairs of vertex kinds.
The inverse of \lean{fr} is the isomorphism, and it sends the literals and
the anchors to themselves by construction.

\paragraph{From minimum gadgets to all valid gadgets.}
\lean{minimum\_unique} is stated for every valid gadget. Its proof takes the
least number $n$ such that some valid gadget on a finite vertex type has $n$
auxiliary vertices. A gadget that attains $n$ is minimum, so $n=6$ by
\lean{minimum\_is\_trident}. Hence every valid gadget has at least six
auxiliary vertices, and at least twelve vertices by the count of
Section~\ref{sec:conventions}. A valid gadget with exactly six is minimum
itself, and \lean{minimum\_is\_trident} applies to it. This is the remark at
the beginning of the paper's proof.

\paragraph{The trident.} \lean{trident\_is\_valid\_gadget} is proved in
\lean{proof/TridentGadget.lean}: the action permutes the indices of $x_i$,
$p_i$ and $u_i$; for an input with a true literal an explicit coloring is
proper; and under $(\F,\F,\F)$ every $p_i$ has color $\B$, which leaves the
upper triangle two colors. This is Theorem~4 of the paper.

% ---------------------------------------------------------------------------
\section{Remarks from the formalization}\label{sec:remarks}

\begin{enumerate}
\item \emph{No gap was found in the proof of Theorem~16.} Every step of the
paper's argument is formalized, in the paper's order, and no case appeared
that the paper does not treat. The Lean proofs differ from the text in the
details of items 4--8.
\item \emph{Two hypotheses are not used.} The proof of Lemma~13 does not use
that the attachment of the matched base is not crossed: the base is recolored
with $\B$ or $\T$ and the literals with $\F$, so every edge between them is
proper. The proof of Lemma~15 does not use that the orbit is free, only that
it is the sole auxiliary orbit. Both hypotheses are kept in the Lean
statements, with an underscore in their names.
\item \emph{Lemma~13 covers orbits of size 2.} The proof uses a singleton
only through the even permutations, which fix it. The same argument gives
\begin{code}
theorem lemma13_of_even_fixed (P : g.Layer) (hm : P.matched) (hiso : ¬ P.triangle)
    (hfix : ∀ w, g.IsAux w → (∀ i, w ≠ P.pt i) →
      ∀ σ, Equiv.Perm.sign σ = 1 → g.ρ σ w = w) : ¬ g.Valid
\end{code}
for other auxiliary vertices in orbits of size 1 or 2, in any number. With
this form the ``remaining budget'' step needs neither Lemma~7 nor the count
that excludes a second size-2 orbit. The package follows the paper: it
proves Lemma~7 and applies Lemma~13 as stated.
\item \emph{Lemma~14} is stated for two disjoint size-3 orbits of a minimum
valid gadget; that there are no other auxiliary vertices is a consequence of
the budget (Section~\ref{sec:elimination}).
\item \emph{Corollary~9(ii)} is applied where it is needed: in Lemma~14, in
Lemma~15, and for the anchors of the branch layer in the triangle analysis.
The proof of Theorem~16 has no separate normalization step.
\item \emph{The case $d=0$} is proved by deleting the second layer and
applying Lemma~13 to the smaller gadget.
\item \emph{Claim~1} is proved by counting permutations
(Section~\ref{sec:orbits}). Theorem~16 does not use it, and the proof of
Lemma~15 uses the literal neighbors of the two ends of an internal edge in
place of their fibers.
\item \emph{Claim~2} needs the standing assumption, which the Lean statement
names as a hypothesis.
\item \emph{The remark of Lemma~11} on arbitrary clause gadgets is proved
for the graph of Definition~1 without the action
(\lean{ClauseGadget.lemma11}). Its argument needs the standing assumption
for the edges to $\GF$, and \lean{ClauseGadget.noLitAnchor} proves the
assumption without symmetry (Section~\ref{sec:principles}).
\item \emph{Details that the paper leaves to the reader.} The merge of
Corollary~9(ii) and the deletion in the case $d=0$ are one construction
(\lean{Removal}). The phrase ``a minimum valid gadget'' becomes
\lean{IsMinimum}, and the passage to all valid gadgets is the argument with
the least number of auxiliary vertices in Section~\ref{sec:thm16}.
\end{enumerate}

% ---------------------------------------------------------------------------
\section{Trusted base}\label{sec:trust}

\begin{itemize}
\item Lean's kernel (Lean v4.33.0) and the standard axioms \lean{propext},
\lean{Classical.choice} and \lean{Quot.sound}. The build prints the axioms of
the two theorems, and these three are all of them. The end of
\lean{Theorem16.lean} also states them as expected output
(\lean{\#guard\_msgs}), so the build fails if a theorem depends on any other
axiom.
\item Mathlib v4.33.0.
\end{itemize}
There is no \lean{native\_decide} and no \lean{sorry}. The tactic
\lean{decide} is evaluated by the kernel. The package calls it
84 times, for small concrete facts only:
\begin{center}
\begin{tabular}{@{}lrp{0.62\textwidth}@{}}
\toprule
file & calls & what is decided\\
\midrule
\lean{OrbitStructure.lean} & 2 & two facts about three indices\\
\lean{TridentGadget.lean} & 5 & the explicit colorings of the trident are proper; two facts about three colors; the trident has twelve vertices and six auxiliary vertices\\
\lean{Basics.lean} & 5 & four facts about three colors; the value of an input at an index\\
\lean{Lemma6.lean} & 7 & three facts about three indices; four products of signs\\
\lean{Lemma7.lean} & 8 & products of signs\\
\lean{Lemma8.lean} & 2 & the swap of $\T$ and $\F$ on the three colors and on the color of a literal\\
\lean{Corollary9.lean} & 1 & a fact about three colors\\
\lean{Lemma10.lean} & 1 & an inequality of two colors\\
\lean{ClauseGadget.lean} & 3 & an inequality of indices; two inequalities of colors\\
\lean{Corollary12.lean} & 3 & inequalities of indices\\
\lean{Lemma13.lean} & 6 & values of the input $(\F,\F,\T)$; inequalities of colors\\
\lean{TwoLayers.lean} & 2 & inequalities of indices and of colors\\
\lean{Lemma14.lean} & 29 & values of the input $(\F,\F,\T)$; inequalities of indices and of colors\\
\lean{TriangleAnalysis.lean} & 4 & inequalities of indices\\
\lean{Lemma15.lean} & 6 & inequalities of indices and of colors\\
\bottomrule
\end{tabular}
\end{center}
No call enumerates gadgets or permutations. A color and an index are
elements of \lean{Fin 3}, a sign is one of the two units of the integers,
and a truth value is a \lean{Bool}. The two largest checks are in
\lean{TridentGadget.lean}: that the explicit coloring of the trident is
proper for each of the eight inputs and each choice of a true literal, and
the two counts of its vertices.

% ---------------------------------------------------------------------------
\section{What is not formalized}\label{sec:gaps}

\begin{itemize}
\item The remark after Theorem~16 that the isomorphism onto the trident is
unique and commutes with the two $S_3$-actions. \lean{minimum\_unique} gives
the existence of the isomorphism only.
\item The exhaustive enumeration of the paper's Section~IV-D and the counts
of its Table~1. The enumeration program is in the parent folder, and the
other two Lean packages of the repository repeat the search inside Lean.
\item The figures of Section~IV and the table that assigns a lemma to each
orbit composition. They summarize the argument and are not statements of
their own.
\item The other results of the paper, except Theorem~4.
\end{itemize}

% ---------------------------------------------------------------------------
\section{Checks}\label{sec:checks}

\begin{itemize}
\item \lean{lake build} on Lean v4.33.0 with Mathlib v4.33.0: exit code 0,
1367 jobs, no warning; the package's own 20 files take
about 12 seconds on the build machine once Mathlib's compiled files are in
place. The build prints the axioms of the two theorems and checks them
against the three standard axioms.
\item \lean{lake env lean checks/ProofCone.lean} follows types, definition
bodies and proofs from the two theorems and from the statements of the paper
that the theorems do not use as stated: \lean{Gadget.lemma6},
\lean{Gadget.fiber\_card}, and \lean{ClauseGadget.lemma11} with
\lean{Gadget.toClauseGadget\_spec}. Of the 314 declarations of the
package, 252 are used. One declaration written in the source is not
used: \lean{exists\_orbit\_labeling} in \lean{OrbitStructure.lean}, which the
compact package uses. The other 61 are generated by Lean; among them
are the projections of two fields of \lean{ClauseGadget} that Lemma~11 does
not read.
\item \lean{Definition1.lean} and \lean{OrbitStructure.lean} are identical to
the files of the compact package, and \lean{TridentGadget.lean} differs in
its import line only.
\end{itemize}

% ---------------------------------------------------------------------------
\section{Files}\label{sec:files}

The folder \lean{proof/} has 20 files with 3,560 lines.
17 of them, 3,145 lines, are specific to this package.

\begin{center}
\begin{tabular}{@{}lrl@{}}
\toprule
file & lines & content\\
\midrule
\lean{Definition1.lean} & 97 & Definition 1, validity, the trident graph (shared)\\
\lean{OrbitStructure.lean} & 173 & Lemma 6: the subgroups of $S_3$ and the orbit types (shared)\\
\lean{TridentGadget.lean} & 145 & the trident as a gadget, Theorem 4 (shared)\\
\lean{Basics.lean} & 367 & anchors, colorings, inputs, the standing assumption, counting\\
\lean{Lemma6.lean} & 351 & Lemma 6 for gadgets; layers and their attachments\\
\lean{Definition2.lean} & 76 & Definition 2 and Claim 1\\
\lean{Lemma7.lean} & 210 & Lemma 7\\
\lean{Lemma8.lean} & 97 & Lemma 8\\
\lean{Removal.lean} & 211 & removing auxiliary vertices from a gadget\\
\lean{Corollary9.lean} & 141 & Corollary 9; minimum gadgets\\
\lean{Lemma10.lean} & 79 & Lemma 10 and Claim 2\\
\lean{Lemma11.lean} & 65 & Lemma 11\\
\lean{ClauseGadget.lean} & 189 & Lemma 11 for arbitrary clause gadgets\\
\lean{Corollary12.lean} & 46 & Corollary 12\\
\lean{Lemma13.lean} & 143 & Lemma 13\\
\lean{TwoLayers.lean} & 254 & gadgets with two size-3 orbits\\
\lean{Lemma14.lean} & 174 & Lemma 14\\
\lean{TriangleAnalysis.lean} & 216 & the triangle analysis\\
\lean{Lemma15.lean} & 156 & Lemma 15\\
\lean{Theorem16.lean} & 370 & Theorem 16\\
\bottomrule
\end{tabular}
\end{center}

\end{document}
